% Laget automatisk av elevark/lag_elevark.py fra sim/relativistisk-bevegelsesmengde.html. Endringer her blir overskrevet.
\documentclass[11pt]{article}
\usepackage{elevark}
\begin{document}
\elevarkhode{Fysikk 2 · Spesiell relativitetsteori}{Konstant kraft og lysfarten}{https://fysikklab.no/sim/relativistisk-bevegelsesmengde.html}{Simuleringen}

Et proton skyves med en konstant kraft. Med $p = mv$ og impulsloven $F \cdot t = Δp$ skulle farten øke jevnt og passere lysfarten etter en stund. Det skjer aldri. Her sammenligner du den klassiske og den relativistiske beskrivelsen.

\section*{Oppgaver}

Start med et proton og $F = 1{,}0 \cdot 10^{-12}\,\mathrm{N}$. $m_{\mathrm{p}} = 1{,}67 \cdot 10^{-27}\,\mathrm{kg}$ og $c = 3{,}00 \cdot 10^{8}\,\mathrm{m/s}$.

\begin{oppgaveliste}
\oppgave \fase{gjett}{Gjett}Regn ut klassisk hvor lang tid det tar før protonet når lysfarten: $F \cdot t = mc$. Hva tror du skjer i virkeligheten? Kjør simuleringen og se på fartsgrafen.
\svar
\oppgave \fase{test}{Regn}Når $t = mc/F$, er $p = mc$. Bruk $p = γmv$ til å vise at farten da er $v = c/\sqrt{2}$. Hva er $γ$? Sjekk med simuleringen.
\svar
\oppgave \fase{forklar}{Forklar}Den samme impulsen gir den samme bevegelsesmengden i begge beskrivelsene. Hvorfor blir farten likevel forskjellig? Hva er det kraften «bruker» bevegelsesmengden til når farten er nær $c$?
\svar
\oppgave \fase{test}{Undersøk}Velg grafen for bevegelsesmengde mot fart. Kraft 2 har en tommelfingerregel om at vi kan regne klassisk under 10~\% av lysfarten. Hvor stor prosentvis feil gjør vi med $p = mv$ ved $0{,}10c$ og ved $0{,}50c$?
\svar
\oppgave \fase{test}{Regn}Velg romskipet på 50 tonn med en motor som skyver med 1,0~MN. Hvor lang tid tar det før romskipet har $0{,}99c$? Sammenlign med tiden det klassisk ville tatt å nå $c$.
\svar
\oppgave \fase{forklar}{Forklar}Velg grafen for kinetisk energi mot fart. Hvorfor trengs det mye mer energi for å øke farten fra $0{,}90c$ til $0{,}9991c$ enn fra $0{,}70c$ til $0{,}90c$ (eksempel 4-4 i Kraft 2)? Hvorfor kan ingen partikkel med masse nå lysfarten?
\svar
\oppgave \fase{utfordring}{Utfordring}Vis at $v = pc / \sqrt{p^{2} + m^{2}c^{2}}$ (Kraft 2-oppgave 4.51). Bruk uttrykket til å forklare at fartsgrafen flater ut mot $c$, og at den ligner $v = p/m$ når $p$ er liten.
\svar
\end{oppgaveliste}

\end{document}
